% 03_excellence.tex - Excelencia, innovacion y ventaja frente a competidores.
\section{Excelencia, innovacion y posicion frente al estado del arte}
% TODO: esta seccion suele tener un peso alto en la evaluacion (ver el peso
% real en context/convocatoria_checklist.md). Todo lo afirmado debe poder
% defenderse con evidencia registrada (src/references.bib) o con capacidades
% verificables del consorcio (data/consortium.yaml); si no, \todo{CITA
% REQUERIDA} o \todo{DATO PENDIENTE}.

\subsection{Innovacion tecnologica y metodologica}
% TODO: que aporta el proyecto que el estado del arte no ofrece. Evitar
% adjetivos promocionales (revolucionario, disruptivo, sin precedentes):
% cuantificar la mejora cuando exista el dato.

\subsection{Ventaja frente a competidores y alternativas}
% TODO: comparacion honesta con las alternativas existentes (academicas y del
% sector). Puede apoyarse en la evidencia de rag/past_projects/ re-verificada
% contra el SSOT (seccion 8 del marco) y en citas de src/references.bib.

\subsection{Mas alla del estado del arte}
% TODO: ambicion del proyecto y resultados que sobreviviran al propio
% proyecto (conocimiento nuevo, metodologias, patrones reutilizables).
